\documentclass[a4paper,12pt]{article}
\usepackage[T2A]{fontenc}
\usepackage[utf8]{inputenc}
\usepackage[russian]{babel}
\usepackage{amsmath,amssymb,mathtools}
\usepackage{helvet}
\renewcommand{\familydefault}{\sfdefault}
\usepackage{geometry}
\geometry{left=22mm,right=22mm,top=20mm,bottom=22mm,headheight=23pt}
\usepackage{microtype}
\usepackage{icomma}
\usepackage{xcolor}
\usepackage{tikz}
\usetikzlibrary{arrows.meta,positioning,shapes.geometric,calc}
\usepackage{tkz-euclide}
\usepackage{pgfplots}
\pgfplotsset{compat=1.18}
\usepackage{booktabs,tabularx}
\usepackage{enumitem}
\usepackage{hyperref}
\usepackage{parskip}
\usepackage{fancyhdr}
\usepackage{setspace}
\onehalfspacing
\definecolor{accent}{HTML}{008080}
\definecolor{darkaccent}{HTML}{004D4D}
\definecolor{textgray}{HTML}{333333}
\definecolor{lightline}{HTML}{C8DCDD}
\hypersetup{colorlinks=true,linkcolor=darkaccent,pdfauthor={Лёвин Артём Александрович},pdftitle={Чек-лист: графики, уравнения и геометрия — Софья Кальней}}
\setlength{\parindent}{0pt}
\setlength{\parskip}{5pt}
\setlist[enumerate]{leftmargin=*,itemsep=7pt,topsep=4pt}
\pagestyle{fancy}
\fancyhf{}
\fancyhead[L]{\footnotesize\color{darkaccent} Математика · повторение}
\fancyhead[R]{\footnotesize\color{textgray} Софья Кальней}
\fancyfoot[C]{\footnotesize\thepage}
\renewcommand{\headrulewidth}{0.3pt}
\renewcommand{\headrule}{\hbox to\headwidth{\color{lightline}\leaders\hrule height \headrulewidth\hfill}}
\newcommand{\key}[1]{\textcolor{darkaccent}{\textbf{#1}}}
\newcommand{\checkline}[1]{\par\noindent\textcolor{accent}{$\square$}\hspace{5pt}#1\par}
\newcommand{\ex}[1]{\par\smallskip\textcolor{darkaccent}{\textbf{Пример.}} #1\par}
\newcommand{\warn}[1]{\par\textbf{Важно.} #1\par}
\newcommand{\topic}[1]{\vspace{5pt}{\color{darkaccent}\large\bfseries #1}\par\vspace{3pt}}
\begin{document}
\thispagestyle{empty}
\begin{center}
\vspace*{28mm}
{\large\color{darkaccent} МАТЕМАТИКА}\par
\vspace{18mm}
{\Huge\color{textgray} Чек-лист}\par
\vspace{11mm}
{\LARGE\color{accent} Графики, уравнения\par и геометрия}\par
\vspace{13mm}
{\large Подготовка к ОГЭ по математике}\par
\vfill
{\large Софья Кальней}\par
\vspace{13mm}
{15 ноября 2026 г.}\par
\vspace{18mm}
{\small Лёвин Артём Александрович}\par
{\small эксклюзивно для Софьи Кальней}\par
\vspace{20mm}
\begin{tikzpicture}
\draw[accent,thin] (0,0) .. controls (2.3,0.85) and (3.4,-0.6) .. (5.1,0)
 .. controls (7.0,0.85) and (8.1,-0.6) .. (9.8,0);
\end{tikzpicture}
\end{center}
\clearpage
\section*{1. Задачи на движение по реке}
\checkline{Отличаю собственную скорость судна от скорости течения.}
\checkline{Правильно составляю таблицу «скорость — время — расстояние».}
\checkline{Учитываю, на сколько различаются времена движения.}

Пусть $v$ --- собственная скорость судна, $u$ --- скорость течения, $s$ --- расстояние в одну сторону.
\[
\begin{array}{c|c|c}
\text{Направление}&\text{Скорость}&\text{Время}\\ \hline
\text{По течению}&v+u&\dfrac{s}{v+u}\\[5pt]
\text{Против течения}&v-u&\dfrac{s}{v-u}
\end{array}
\]
\key{Ограничение:} $v>u\ge 0$; тогда обе скорости положительны.

\ex{Судно прошло по $192$ км в каждом направлении. Скорость течения $4$ км/ч, а против течения путь занял на $4$ часа больше. Найдём $v$.}
\[
\frac{192}{v-4}-\frac{192}{v+4}=4,\qquad v>4.
\]
\[
\frac{192(v+4)-192(v-4)}{v^2-16}=4
\ \Longrightarrow\ \frac{1536}{v^2-16}=4.
\]
\[
v^2-16=384,\qquad v^2=400,\qquad \boxed{v=20\ \text{км/ч}}.
\]
Отрицательный корень исключён условием $v>4$. Проверка: $192/16-192/24=12-8=4$ часа.

\section*{2. Как извлекать большие квадратные корни}
\topic{Способ 1. Оценка и последняя цифра}
\checkline{Нахожу два соседних десятка или сотни, между квадратами которых лежит число.}
\checkline{Проверяю возможные последние цифры корня.}
\checkline{Обязательно возвожу найденный кандидат в квадрат.}

\ex{$\sqrt{5184}$. Так как $70^2=4900$, $80^2=6400$, корень между $70$ и $80$. Последняя цифра квадрата равна $4$, поэтому возможны окончания $2$ и $8$. Проверяем: $72^2=5184$, значит $\sqrt{5184}=\boxed{72}$.}

\ex{$\sqrt{9409}$. Число расположено между $90^2=8100$ и $100^2=10000$. Квадрат оканчивается на $9$, значит последняя цифра корня --- $3$ или $7$. Проверка: $97^2=9409$, поэтому $\sqrt{9409}=\boxed{97}$.}

\warn{Последняя цифра только сокращает число вариантов. Окончательный ответ подтверждается проверкой.}
\clearpage
\topic{Способ 2. Найти корень дискриминанта рядом с известным числом}
Если $D=b^2-4ac$ и $4ac$ сравнительно невелико, квадратный корень из $D$ может оказаться близок к $|b|$. Это предположение требуется проверить точно.

\ex{Решим $60x^2-1802x+60=0$.}
\[
\begin{aligned}
D&=1802^2-4\cdot60\cdot60\\
 &=3\,247\,204-14\,400=3\,232\,804.
\end{aligned}
\]
Предполагаем, что $\sqrt D=1798$. Проверим разностью квадратов:
\[
1802^2-1798^2=(1802-1798)(1802+1798)
=4\cdot3600=14\,400.
\]
Значит, действительно $D=1798^2$, и
\[
x=\frac{1802\pm1798}{120},\qquad
\boxed{x_1=30,\quad x_2=\frac1{30}}.
\]
\topic{Способ 3. Разложение на множители}
Если $N=p_1^{2k_1}\cdots p_r^{2k_r}$, то $\sqrt N=p_1^{k_1}\cdots p_r^{k_r}$.
Способ универсален для полных квадратов, однако для больших чисел разложение бывает трудоёмким.

\section*{3. Квадратное уравнение и дискриминант}
\checkline{Использую $D=b^2-4ac$ и $x=\dfrac{-b\pm\sqrt D}{2a}$ при $a\ne0$.}
\checkline{Умею оценить $\sqrt D$ и проверить кандидат точно.}
\checkline{Проверяю найденные корни подстановкой или по теореме Виета.}

Для уравнения $ax^2+bx+c=0$: если $D<0$, действительных корней нет; если $D=0$, один корень; если $D>0$, два различных корня.
\warn{На графике параболы отсутствие пересечений с осью $Ox$ соответствует отрицательному дискриминанту.}

\clearpage
\section*{4. Кусочные функции: строим каждую ветвь на её области}
\checkline{Для каждой формулы сначала выписываю разрешённые значения $x$.}
\checkline{Для параболы нахожу вершину, направление ветвей и несколько опорных точек.}
\checkline{При строгом неравенстве ставлю пустую точку; при нестрогом --- закрашенную.}
\checkline{Чтобы решить $f(x)=m$, считаю пересечения графика с прямой $y=m$.}

\ex{Рассмотрим функцию, разобранную на занятии:}
\[
f(x)=\begin{cases}
-x^2+8x-17,&x\ge2,\\
-x-2,&x<2.
\end{cases}
\]
Для параболы
\[
-x^2+8x-17=-(x-4)^2-1.
\]
Поэтому вершина $(4;-1)$; ветви направлены вниз. При $x=2$ имеем $f(2)=-5$, и эта точка \emph{закрашена}. На прямой $y=-x-2$ значение в $x=2$ было бы $-4$, но $x=2$ для этой ветви запрещено: точка $(2;-4)$ \emph{пустая}.

\begin{center}
\begin{tikzpicture}
\begin{axis}[
width=0.88\linewidth,height=7.5cm,
xmin=-2.6,xmax=8.1,ymin=-8.4,ymax=2.6,
axis lines=middle,axis equal image,
xtick={-2,0,2,4,6,8},ytick={-8,-6,-5,-4,-2,-1,0,2},
minor tick num=0,grid=major,grid style={lightline!55},
axis line style={->,>=Stealth},
xlabel={$x$},ylabel={$y$},
legend pos=north east,legend style={font=\footnotesize,draw=none,fill=white},
clip=true]
\addplot[blue!60!black,thick,domain=-2.6:1.999,samples=120]{-x-2};
\addlegendentry{$-x-2$, $x<2$}
\addplot[accent,thick,domain=2:7.2,samples=220]{-x^2+8*x-17};
\addlegendentry{$-(x-4)^2-1$, $x\ge2$}
\addplot[only marks,mark=*,mark size=2.8pt,accent] coordinates {(2,-5)};
\addplot[only marks,mark=o,mark size=3.3pt,very thick,blue!60!black] coordinates {(2,-4)};
\addplot[only marks,mark=*,mark size=2.3pt,accent] coordinates {(4,-1)};
\end{axis}
\end{tikzpicture}
\end{center}

\textbf{Важный вывод:} две формулы могут давать разные значения на границе. Соединять ветви искусственно нельзя.

\clearpage
\section*{5. Сколько решений у уравнения $f(x)=m$?}
Горизонтальная прямая $y=m$ показывает число решений: одному пересечению соответствует один допустимый $x$.

В нашем примере прямая $y=-x-2$ при $x<2$ принимает только значения $y>-4$.
Парабола на области $x\ge2$ достигает максимума $-1$ и проходит через $(2;-5)$.

\begin{center}
\begin{tabularx}{\linewidth}{@{}p{35mm}X@{}}
\toprule
\textbf{Уровень $m$} & \textbf{Количество пересечений}\\
\midrule
$m<-5$ & Одна точка параболы, у прямой --- ни одной.\\
$m=-5$ & Две точки параболы (включая $x=2$).\\
$-5<m\le-4$ & Две точки параболы, точек прямой нет.\\
$-4<m<-1$ & Две точки параболы и одна точка прямой.\\
$m=-1$ & Вершина параболы и одна точка прямой.\\
$m>-1$ & Одна точка прямой, точек параболы нет.\\
\bottomrule
\end{tabularx}
\end{center}

Следовательно, \key{ровно два решения} имеются при
\[
\boxed{m\in[-5;-4]\cup\{-1\}}.
\]
\warn{Значение $m=-4$ входит в ответ: точка прямой в $x=2$ пустая, однако у параболы есть две допустимые точки. При $m=-1$ две точки снова получаются за счёт разных ветвей.}

\clearpage
\begin{center}
\vspace*{8mm}
{\Large\color{darkaccent}\bfseries Алгоритм: построение кусочного графика}
\end{center}
\vspace{9mm}
\begin{center}
\begin{tikzpicture}[
 node distance=10mm and 13mm,
 terminal/.style={ellipse,draw=darkaccent,thick,minimum width=49mm,minimum height=11mm,align=center,font=\small},
 process/.style={rectangle,rounded corners=2pt,draw=darkaccent,thick,text width=56mm,minimum height=12mm,align=center,font=\small,inner sep=5pt},
 branch/.style={diamond,aspect=2.2,draw=accent,thick,text width=38mm,align=center,font=\small,inner sep=3pt},
 side/.style={rectangle,rounded corners=2pt,draw=darkaccent,text width=49mm,minimum height=19mm,align=center,font=\small,inner sep=5pt},
 arrow/.style={-{Stealth[length=2.7mm]},thick,darkaccent}
]
\node[terminal] (start) {Начало};
\node[process,below=of start] (rules) {Запишите формулы и области определения каждой ветви};
\node[branch,below=13mm of rules] (decision) {Какую ветвь строим?};
\node[side,below left=22mm and 12mm of decision] (left) {Прямая: две точки, только при $x<2$};
\node[side,below right=22mm and 12mm of decision] (right) {Парабола: вершина и опорные точки, $x\ge2$};
\node[process,below=42mm of decision] (dots) {Отметьте пустую и закрашенную точки при $x=2$};
\node[process,below=of dots] (horizontal) {Для $f(x)=m$ проведите $y=m$ и посчитайте пересечения};
\node[terminal,below=of horizontal] (finish) {Проверен ответ};
\draw[arrow] (start)--(rules);
\draw[arrow] (rules)--(decision);
\draw[arrow] (decision.west)-| node[pos=.18,above,font=\footnotesize] {$x<2$} (left.north);
\draw[arrow] (decision.east)-| node[pos=.18,above,font=\footnotesize] {$x\ge2$} (right.north);
\draw[arrow] (left.south) |- ($(dots.north)+(0,4mm)$) -- (dots.north);
\draw[arrow] (right.south) |- ($(dots.north)+(0,4mm)$);
\draw[arrow] (dots)--(horizontal);
\draw[arrow] (horizontal)--(finish);
\end{tikzpicture}
\end{center}
\clearpage
\section*{6. Геометрия: биссектриса в параллелограмме}
\checkline{Использую параллельность противоположных сторон параллелограмма.}
\checkline{Доказываю равнобедренность треугольника равенством углов.}
\checkline{Для периметра складываю длины двух соседних сторон и умножаю на $2$.}

\ex{В параллелограмме $ABCD$ биссектриса угла $DAB$ пересекает $BC$ в точке $K$. Известно: $BK=12$, $CK=16$. Найдём периметр.}
\begin{center}
\begin{tikzpicture}[scale=.9]
\tkzDefPoint(0,0){A}
\tkzDefPoint(1.2,2.078){B}
\tkzDefPoint(6.8,2.078){C}
\tkzDefPoint(5.6,0){D}
\tkzDefPoint(3.6,2.078){K}
\tkzDrawPolygon[thick,color=darkaccent](A,B,C,D)
\tkzDrawSegment[thick,color=accent](A,K)
\tkzDrawPoints[fill=darkaccent,size=2](A,B,C,D,K)
\tkzLabelPoint[below left](A){$A$}
\tkzLabelPoint[above left](B){$B$}
\tkzLabelPoint[above right](C){$C$}
\tkzLabelPoint[below right](D){$D$}
\tkzLabelPoint[above](K){$K$}
\node at (2.45,2.48) {$12$};
\node at (5.25,2.48) {$16$};
\end{tikzpicture}
\end{center}
Поскольку $AD\parallel BC$, получаем $\angle KAD=\angle AKB$ (накрест лежащие). По свойству биссектрисы $\angle BAK=\angle KAD$, поэтому
\[
\angle BAK=\angle AKB\quad\Longrightarrow\quad AB=BK=12.
\]
Также $BC=BK+KC=12+16=28$. Тогда
\[
\boxed{P=2(AB+BC)=2(12+28)=80}.
\]
\warn{Факт $AB=BK$ нужно обосновать через углы. Он следует из данных о биссектрисе и параллельных сторонах.}

\section*{7. Площади треугольников с общей высотой}
Если треугольники $ABK$ и $ABD$ имеют общую высоту из $A$ к прямой $BD$, то площади относятся как основания:
\[
\boxed{\frac{S_{ABK}}{S_{ABD}}=\frac{BK}{BD}}.
\]
\ex{Если $BK=7$, $BD=20$, $S_{ABD}=60$, то $S_{ABK}=60\cdot\frac7{20}=21$.}

\clearpage
\section*{8. Площадь четырёхугольника через диагонали}
\checkline{Знаю формулу площади выпуклого четырёхугольника через диагонали.}
\checkline{Умею найти площадь треугольника по двум сторонам и синусу угла.}
\checkline{Умею исключить общий множитель $\sin\varphi$ из двух выражений площадей.}

Если диагонали $AC$ и $BD$ выпуклого четырёхугольника пересекаются в $O$ под углом $\varphi$, то
\[
\boxed{S_{ABCD}=\frac12 AC\cdot BD\cdot\sin\varphi},
\qquad AC=AO+OC.
\]
Для треугольника $ABD$ также верно
\[
S_{ABD}=\frac12 AO\cdot BD\cdot\sin\varphi.
\]
Эти формулы относятся к обычной невырожденной конфигурации, где $0<\varphi<180^\circ$ и $\sin\varphi>0$.

\begin{center}
\begin{tikzpicture}[scale=.88]
\tkzDefPoint(0,0){A}
\tkzDefPoint(.6,3.3){B}
\tkzDefPoint(5.6,3.5){C}
\tkzDefPoint(5.1,-1.2){D}
\tkzDefPoint(2.4,1.5){O}
\tkzDefPoint(2.7,1.2){K}
\tkzDrawPolygon[thick,color=darkaccent](A,B,C,D)
\tkzDrawSegment[thin,color=textgray](A,C)
\tkzDrawSegment[thin,color=textgray](B,D)
\tkzDrawSegments[thick,color=accent](A,B B,K K,A)
\tkzDrawPoints[fill=darkaccent,size=2](A,B,C,D,O,K)
\tkzLabelPoint[below left](A){$A$}
\tkzLabelPoint[above left](B){$B$}
\tkzLabelPoint[above right](C){$C$}
\tkzLabelPoint[below right](D){$D$}
\tkzLabelPoint[above right](O){$O$}
\tkzLabelPoint[below](K){$K$}
\node at (1.2,.85) {$12$};
\node at (4.25,2.75) {$16$};
\end{tikzpicture}
\end{center}

\ex{Дано $AO=12$, $OC=16$, $BD=18$, $K\in BD$ и $S_{ABK}=\frac15S_{ABCD}$. Найдём $KD$.}
\[
S_{ABCD}=\frac12\cdot(12+16)\cdot18\sin\varphi
=252\sin\varphi.
\]
\[
S_{ABK}=\frac15\cdot252\sin\varphi
=50{,}4\sin\varphi.
\]
С другой стороны, по общей высоте
\[
\begin{aligned}
S_{ABK}&=\frac{BK}{BD}S_{ABD}
=\frac{BK}{18}\cdot\frac12\cdot12\cdot18\sin\varphi\\
&=6BK\sin\varphi.
\end{aligned}
\]
Следовательно, $6BK\sin\varphi=50{,}4\sin\varphi$.
Поскольку $\sin\varphi>0$, то $BK=8{,}4$, а
\[
\boxed{KD=BD-BK=18-8{,}4=9{,}6}.
\]
\warn{Точки $O,B,K$ находятся на одной прямой $BD$, поэтому площадь $OBK$ равна нулю. В содержательном условии используется площадь треугольника $ABK$.}
\clearpage
\section*{Тренировочный блок · 10 задач}
\textcolor{darkaccent}{\textbf{Средний уровень}}
\begin{enumerate}
\item Вычислите устно, затем проверьте возведением в квадрат: $\sqrt{6084}$.
\item Найдите $\sqrt{9409}$. Объясните, как определить последнюю цифру корня.
\item Извлеките корень $\sqrt{2\,556\,801}$, используя близость к квадрату числа $1601$.
\end{enumerate}
\textcolor{darkaccent}{\textbf{Уровень выше среднего}}
\begin{enumerate}[resume]
\item Решите уравнение $40x^2-1601x+40=0$.
\item Теплоход проплыл $192$ км по течению и $192$ км против течения. Скорость течения $4$ км/ч; против течения путь занял на $4$ ч больше. Найдите собственную скорость теплохода.
\item Для функции $y=-x^2+6x-10$ найдите вершину параболы, значение $y$ при $x=1$ и выясните, пересекает ли её график ось $Ox$.
\item Постройте график функции
\[
f(x)=\begin{cases}-x^2+6x-10,&x\ge1,\\-x-2,&x<1.\end{cases}
\]
Найдите все $m$, при которых уравнение $f(x)=m$ имеет ровно два различных решения.
\item Биссектриса угла $A$ параллелограмма $ABCD$ пересекает $BC$ в $K$. Найдите периметр параллелограмма, если $BK=9$, $KC=15$.
\end{enumerate}
\textcolor{darkaccent}{\textbf{Сложный уровень}}
\begin{enumerate}[resume]
\item Точка $K$ лежит на стороне $BD$ треугольника $ABD$. Известны $BD=20$, $BK=7$ и $S_{ABD}=60$. Найдите $S_{ABK}$.
\item В выпуклом четырёхугольнике $ABCD$ диагонали пересекаются в точке $O$, $AO=10$, $OC=15$, $BD=20$. Точка $K\in BD$, причём $S_{ABK}=\dfrac14 S_{ABCD}$. Найдите $KD$.
\end{enumerate}
\clearpage
\section*{Ответы к тренировочному блоку}
\renewcommand{\arraystretch}{1.42}
\begin{center}
\begin{tabularx}{\linewidth}{@{}p{15mm}X@{}}
\toprule
\textbf{№} & \textbf{Ответ}\\
\midrule
1 & $78$\\
2 & $97$\\
3 & $1599$\\
4 & $x_1=40$, $x_2=\dfrac1{40}$\\
5 & $20$ км/ч\\
6 & Вершина $(3;-1)$; при $x=1$: $y=-5$; пересечений с $Ox$ нет.\\
7 & $m\in[-5;-3]\cup\{-1\}$\\
8 & $66$\\
9 & $21$\\
10 & $7{,}5$\\
\bottomrule
\end{tabularx}
\end{center}
\vspace{12mm}
\section*{Итоговая памятка}
\begin{tabularx}{\linewidth}{@{}p{53mm}X@{}}
\toprule
\textbf{Сигнал в задании} & \textbf{Как действовать}\\
\midrule
Движение туда и обратно & Запишите два времени как расстояние, делённое на скорость, затем составьте разность.\\
Большой квадратный корень & Оцените величину, проверьте последнюю цифру, затем подтвердите квадратом.\\
Большой дискриминант & Сравните его с квадратом близкого целого и проверьте разность квадратов.\\
Кусочная функция & Стройте каждую формулу только на её области и различайте пустые и закрашенные точки.\\
Ровно два решения $f(x)=m$ & Ищите уровни, на которых горизонтальная прямая имеет два пересечения.\\
Биссектриса в параллелограмме & Докажите равенство углов, затем равенство сторон треугольника.\\
Отношение площадей & Для общей высоты сравните основания; для диагоналей используйте синус угла.\\
\bottomrule
\end{tabularx}
\vfill
{\footnotesize\color{textgray}Лёвин Артём Александрович · эксклюзивно для Софьи Кальней}
\end{document}
